\clearpage
\section{Follow-up Campaign Identities}
\label{app:followupconfigs}
The follow-ups retain the Qwen3-8B/0.6B Q8\_0 pair and upstream llama.cpp b10839, but not one common frontend or reference. Table~\ref{tab:followupconfigs} identifies the new campaigns. Original RTX direct-ID qualification used Windows build 26200; new RTX H1, latency and H2 use 26300. Their driver version is 616.92, versus historical timing's 616.56; version equality does not authenticate original driver bytes. Recorded byte identities cover models, frontend, DLLs, compiler and observed driver modules on the experimental host; omitted assets cannot be rehashed by an offline reader.

\begin{table*}[t]
\caption{Campaign-specific follow-up configuration. RTX always denotes RTX 5070 Ti Laptop. Greedy target sampling uses seed 123, FA off, fresh contexts, and an inclusive EOG/128 committed-ID endpoint. All layers are offloaded: target 37/37 and draft 29/29. The original studies' references remain unchanged.}
\label{tab:followupconfigs}
\centering
\begin{tabular}{>{\raggedright\arraybackslash}p{1.15in}>{\raggedright\arraybackslash}p{2.15in}>{\raggedright\arraybackslash}p{2.45in}}
\toprule
Campaign & Frontend / execution & Reference or paired contrast \\
\midrule
B570 live fork & Instrumented native speculative example; Vulkan0; context 1024, batch 512 & Native UB1 versus state-restored UB128; identical fixed proposals at sampled block; UB128 continuation persists \\
RTX H1 discovery & Original CLI/CUDA DLLs; context 4096 & Target-only UB2 versus batched native live; 24/36 fails selection, reserved 0 \\
RTX H2 V2 & Separate P/R workers and original native CLI; shared numerical DLLs; context 4096, batch 2048 & H2 and R\_ONLY target UB2; original target-only reference UB1; native live target/draft UB2 \\
RTX matched cost & New Clang 21.1.0 timer frontend/core, original hashed CUDA backend DLLs; context 4096 & Same instrumented executable for target-only UB1 and serialized UB1; new 90/90 direct-ID gate \\
\bottomrule
\end{tabular}
\end{table*}

\paragraph{Separate selection and exposure.}
Original B570 H1's 24 reserved tasks and 133-task eligible pool remain as documented above. The live fork study excludes old direct-ID, old reserved and other recorded exposure before deterministic selection using its fixed study salt. The RTX H1 arm excludes 103 exposed or ambiguous tasks from 164, leaving 61; its reserved list is never executed after discovery failure. V2 re-audits accessible scientific files, archives and intervening study records: 114 tasks are excluded, leaving 50. Its selection uses the smallest SHA-256 values of a fixed V2 salt plus full task ID, without outcome selection. The actual 24-task V2 order is /160, /45, /149, /17, /6, /72, /33, /121, /130, /12, /4, /55, /69, /14, /61, /135, /100, /117, /22, /7, /23, /95, /88, /37. These tasks must not be pooled with original B570, the new fork set, the unexecuted H1 reserved list or latency's HumanEval30. Offline selection checks reconstruct retained inventories, not a new scan of unavailable remote directories.

\section{B570 Same-State Fork Endpoints}
The frozen checkpoint is the first native verification batch beginning at at least 64 committed IDs. It is one sampled block per eligible execution, not every generation position. All 72 reserved executions have eligible forks. The live draft's two proposals are recorded before either target decode. Source KV is serialized, restored into the alternate context, serialized again and compared with source-after bytes; consumed length, prefix, positions and physical groups are checked. Native UB1 produces three one-output physical groups, whereas alternate UB128 produces one three-output group at this sampled block.

Sampler identity uses the pinned recursive \texttt{common\_sampler\_clone}, copied chain/seed/accepted-token history and the same-partition R0 control. Six exposed engineering tasks generate 24 fresh R0 executions; original/instrumented primary IDs and proposals agree, and UB1-to-UB1 restored continuation and scores agree. Opaque sampler serialization is unavailable, so these are source-audited clone and empirical controls, not a byte proof for omitted external sampler binaries.

\begin{table}[t]
\caption{B570 reserved fork results: 24 new prompts $\times$ three. Events and repetitions are correlated. Primary immediate decisions remain unchanged; later output differences are a separate secondary endpoint.}
\label{tab:forkendpoints}
\centering
\begin{tabular}{lrr}
\toprule
Endpoint difference & Runs & Prompts \\
\midrule
Raw score records & \FUForkScores & 24 \\
Raw top-1 IDs & 0 & 0 \\
Immediate decisions & \FUForkImmediate & 0 \\
Full continuation IDs & \FUForkLater & \FUForkLaterPrompts \\
\bottomrule
\end{tabular}
\end{table}

The primary comparison follows every reachable decision at a still-common proposal-conditioned prefix in the sampled block until rejection, bonus or endpoint: \FUForkEvents\ paired events. All three raw score rows are retained even when proposals are rejected. Score equality, raw argmax equality and actual accepted/corrected/bonus IDs are different tests. The exported paired records preserve them separately.

After the native continuation ends, the implementation restores the saved draft KV and cloned samplers, replaces target context with the already-verified alternate UB128 context, restores counters and committed prefix, then continues the native loop. Subsequent target blocks remain UB128, and subsequent draft proposals follow alternate accepted history. This continuation policy changes more than a single numerical block over the rest of generation. The 18/72 trajectory differences across six prompts occur after unchanged sampled immediate decisions; they cannot be attributed to an immediate split at that block or claimed as reference restoration. The tested states/policy have score and downstream effects, with no detected immediate decision effect. There is no unique-operation, whole-generation null-effect, RTX or performance inference.

\section{RTX Independent Proposal/Replay Reproduction}
\paragraph{The earlier failed arm.}
H1 target-only UB2 matches 24/36 newly executed exposed batched trajectories. /64, /106, /119 and /128 each fail all three repetitions. The 36/36 gate is unmet, no predictor is selected, no hypothesis or prediction freeze is reached, and reserved runs are zero. H2 was not implemented in that frozen arm. V2 is a new study, not a retroactive execution of the earlier decision tree.

\paragraph{Program independence and shared computation.}
P loads only draft weights; R loads only target weights. H2 orchestrates these processes, first recording each proposal request and then sending it to R. P uses H2's already-produced feedback; it is not a feedback-free draft tape. R computes pending tokens, acceptance, correction, bonus and KV suffix removal. Fresh R\_ONLY reads one sealed tape request at a time and checks its own computed prefix and pending token. Future prediction/reference/live IDs are excluded from its computation inputs. The original native CLI produces its own proposals and has no tape import. Model-loader, compile-time role, allowed RPC fields and process/load records establish application/API isolation, not adversarial OS-wide file access. Original model loaders, samplers, llama.cpp/GGML/CUDA kernels are shared, so there is no independent numerical implementation.

\paragraph{Freeze chain and denominators.}
F0 seals source, effective configuration, input/exposure/selection and qualification before discovery. Discovery completes 108 executions (12 tasks $\times$ three repetitions $\times$ H2, R\_ONLY and native live), with both H2/live and H2/R\_ONLY 36/36. F1 seals the selected program before 216 reserved pre-live executions (reference/H2/R\_ONLY). F2 seals the 72 prediction rows, full tapes and pre-live package; remote acknowledgement and retrieval of the exact sealed bytes precede the first of 72 reserved native live runs. Thus formal four-condition execution totals 288, with 72 comparisons per named endpoint. Git blob/tree membership certifies preserved bytes, while server and host records provide the reported order; neither is a trusted per-process external clock. Sanitized anonymous ordinals cannot independently prove that original absolute chronology.

\begin{table}[t]
\caption{RTX V2 reserved equality. Every row uses 24 distinct prompts $\times$ three repetitions. Prediction equality and reference qualification have different comparators.}
\label{tab:rtxreplay}
\centering
\begin{tabular}{lrr}
\toprule
Comparison & Equal / 72 & Prompts / 24 \\
\midrule
H2 / native live & 72 & 24 \\
R\_ONLY / native live & 72 & 24 \\
H2 / R\_ONLY & 72 & 24 \\
UB1 / native live & 60 & 20 \\
\bottomrule
\end{tabular}
\end{table}

\begin{table}[t]
\caption{RTX V2 reference-divergent cases, stable in all three repetitions. Indices are zero-based generated-ID positions; H2 and R\_ONLY predict the native live value. Reference is this study's UB1 path.}
\label{tab:rtxdivergences}
\centering
\begin{tabular}{lrrr}
\toprule
HumanEval & Index & Reference ID & Live ID \\
\midrule
/17 & 4 & 1406 & 271 \\
/33 & 8 & 14937 & 4210 \\
/14 & 4 & 1406 & 271 \\
/22 & 0 & 1499 & 750 \\
\bottomrule
\end{tabular}
\end{table}

\paragraph{Actual endpoint distribution.}
H2, R\_ONLY and native live each contain 63 earlier-EOG runs on 21 prompts and nine cap runs on three prompts. The UB1 reference has 66 earlier-EOG runs on 22 prompts and six cap runs on two prompts. Each includes the first EOG ID and excludes prompt and uncommitted proposals. This differs from original B570 H1, where every reserved array reaches 128. Repetition stability holds within all four conditions; it is not reference equality. The reference gate remains blocked, and none of the 72/72 prediction rows is a qualified performance result.

\section{Fresh RTX Serialized Cost Estimator and Timer}
Thirty distinct prompts from the latency arm's frozen original HumanEval30 input list each receive three executions of target-only and serialized: 180 new runs, 90 direct-ID pairs. Conditions are randomly ordered within prompt/repetition with seed 2026093001; bootstrap seed is 2026093002.

\begin{samepage}
For condition medians $m_{ic}$ on prompt $i$, compute
\[
\begin{aligned}
r_i&=m_{i,\mathrm{target}}/m_{i,\mathrm{serialized}},\\
G&=\exp\!\left(\frac1{30}\sum_{i=1}^{30}\log r_i\right).
\end{aligned}
\]
\end{samepage}
The 10,000 bootstrap draws resample 30 prompts in original numeric task order using Python Random, with linearly interpolated sorted percentiles at .025/.975. This is neither median of paired ratios nor ratio of 180 total times; repeats are not bootstrap units. Canonical $G=0.6458157742698286$, CI [0.629802205379344, 0.6598843928244866], with 0/30 ratios above one. The independent floating-point replay agrees within $10^{-14}$; full-precision canonical and independently recomputed results remain separately recorded. The rounded manuscript value is \FUCost\ [\FUCostLo, \FUCostHi].

Both conditions use the same freshly built timer executable and pass their own 90/90 direct-ID gate. The native QPC interval begins after synchronized target/draft initial prefill and before per-request sampler creation/reset/seed and generation-state work. Draft inference, target validation, acceptance/correction/bonus, KV state, synchronization, proposal/scheduling and relevant final response/debug-log work through slot release are charged. Loading, initial prefill and teardown are excluded. Persistent engine allocation and one-time draft sampler precede prefill; request sampler work is inside the timer. Endpoint records print after measured end. A second chrono interval cross-checks the native endpoints, but both APIs share Windows' clock oscillator.

The frozen empty-hook probe checks clock-call overhead at its prescribed threshold; it does not show that total committed-ID/debug logging costs less than 1\%. This is instrumented CLI cost, not deployment throughput. Thermal entry and randomized adjacent scheduling reduce drift without locking all clock/session variation. No energy result is estimated. The earlier RTX timing ratio 0.650, Figure 3 proxy cells and qualification-only gates are unchanged. No same-study qualified batched timer endpoint exists, so a full batched-to-serialized speed tradeoff is unmeasured.

\section{Reviewer Evidence and Offline Recalculation}
The anonymous evidence index provides study-separated summaries, full original B570 arrays, RTX H1 discovery arrays, new B570 paired score/decision/continuation records and RTX latency condition times/arrays, with a standard-library verification entry point and derived file hashes. New RTX H2 arrays are included in the index; its larger complete anonymous tape/events capsule is delivered separately with the exact filename and SHA-256 in the index. Original scientific archives stay unchanged. Exported arrays are compared element by element with the source evidence before packaging; tapes and anonymous manifests have their own derived hashes.

Original host paths, user identities, Git author records and exact timestamps are omitted from new compact exports. Public benchmark task IDs remain for mapping. Content/array hashes can still permit linkage to published evidence and are not a guarantee of unlinkability; no identifying repository or local URI is inserted into the PDFs or source ZIP. The old prediction-file digest above is a content locator, not temporal proof. Anonymous original-to-derived mappings and ordinal chronology do not independently certify original remote time or byte identity of sanitized files. The source ZIP rebuilds the two documents; the evidence index/capsule recomputes retained observations without hardware. Neither authenticates omitted weights, runtime, compiler, driver or omitted full logits, executes models, or establishes functional HumanEval pass@$k$. Author upload is a separate step; no anonymous online locator is invented.
